%%%%%%%%%%%%%%%%%%%%%%%%%%%%%%%%%%%%%%%%%%%%%%%%%%%%%%%%%%%%%%%%%%%%%%%%%%%%%%%%%%%%%%%%%%%
%%
%%   gshs-server-intro by gs19007 Junu Kwon.
%%   경기과학고등학교 리눅스 사용자협회
%%   
%%   github.io/junukwon7
%%   junukwon7@gmail.com
%%
%%%%%%%%%%%%%%%%%%%%%%%%%%%%%%%%%%%%%%%%%%%%%%%%%%%%%%%%%%%%%%%%%%%%%%%%%%%%%%%%%%%%%%%%%%%


\RequirePackage[l2tabu, orthodox]{nag} % tells you of any bad LaTeX usage
                                       % must be first thing in class (with the exception of comments)

%% There is one option you should define; oneside or twoside
%% Use twoside for your viva docs (examiners hate long docs they need to carry around)
%% and oneside for the final thing you submit to the library.  Note that margins will
%% change accordingly

\documentclass[twoside]{gshs-server-intro}  % custom gshs


%% **************** (Your) Packages (Start) ******************

% \listfiles % uncomment this to know which packages you are using
              % the list of packages will be in the bottom of the .log file

%% Note that packges may already be loaded from the um (and memoir) classes.
%% Do not add your packages to the template, but rather add them here.

\usepackage{blindtext} %% for some dummy text, remove in your writeup
\usepackage{kotex}
\usepackage{indentfirst}

\usepackage{amsthm}
%% ***************** (Your) Packages (End) *******************

\newtheorem{theorem}{Theorem}[section]
\newtheorem*{remark}{Remark}
\newtheorem{corollary}{Corollary}[theorem]
\newtheorem{lemma}[theorem]{Lemma}
\theoremstyle{definition}
\newtheorem{definition}{Definition}[section]
\renewcommand\qedsymbol{$\blacksquare$}

%% **************** (Your) Data (Start) ******************

\title{Various Methods for \\Calculating $\boldsymbol\zeta$(2)} % 제목
\tagline{Evaluating the value of the infinite sum of $1/{n^2}$\\As known as the Basel problem} % 부제
\author{Junu Kwon} % 이름
\koauthor{19007 권준우}
\department{2021 Calculus I}
\degreedate{\today}
\authorID{}
\cosupervisor{}
\faculty{}
\degree{}
\doctype{}

%% ***************** (Your) Data (End) *******************


%% ******** (Your) Document Settings (Start) *************

% You should have an images directory in every chapX subdir
% NOTE:  Trailing / for subdirs is required.
\graphicspath{{./images/}{./chap1/images/}}   % Paths where to look for images, if defined "images" must always be there as it holds the images in-use by the template.

\makeindex

%% ********* (Your) Document Settings (End) **************

% DOCTOR'S (JP) ORDERS: MAKE SURE TO READ MY TWO BLOG ENTRIES WITH
% CONTENT AND LaTeX TIPS FOR YOUR WRITE-UP.  THESE ARE BASED ON  
% EXAMINER'S FEEDBACK
%
% URLS:
% https://bitsilla.com/blog/2019/03/content-tips-for-your-dissertation-or-project-write-up/
% https://bitsilla.com/blog/2019/01/latex-tips-for-your-dissertation-or-project-write-up/

% end the preamble and start the document

\begin{document}
\frontmatter 
    \maketitle
    %\clearpage
	~\vfill
	\thispagestyle{empty}
	
	\begin{figure}[h]
		\includegraphics[width=4in]{logo_full}
	\end{figure}
	
	\noindent Copyright \textcopyright\the\year\ Junu Kwon \\ % Copyright of the form
	
	\noindent \textsc{https://www.gs.hs.kr/}\\ % URL
	
	\noindent \textit{Compiled \today} % Printing/edition date
	
	%% \noindent Licensed under the Creative Commons Attribution-NonCommercial 3.0 Unported License (the ``License''). You may not use this file except in compliance with the License. You may obtain a copy of the License at \url{http://creativecommons.org/licenses/by-nc/3.0}. Unless required by applicable law or agreed to in writing, software distributed under the License is distributed on an \textsc{``as is'' basis, without warranties or conditions of any kind}, either express or implied. See the License for the specific language governing permissions and limitations under the License.\\ % License information
	\if@openright\cleardoublepage\else\clearpage\fi       
    %\begin{originality}
\end{originality}    
    %\begin{dedication}
\end{dedication}

        % include a dedication.tex file
    %\begin{acknowledgements}
\end{acknowledgements}   % include an acknowledgements.tex file
    \begin{abstract}
The discussion topic in the Math Seminar I class was the examination of  $\sum_{n=1}^\infty 1/n^2$. In the class, we solved it using the Fourier series of $f(x)=x$. However, the teacher said there are numerous proofs out there which uses brilliant techniques so I searched for them and found out that they use Calculus-techniques like integration and Taylor expansion. So, I kept the topic to study about it later.

This report will suggest various ways to prove the identity $\zeta(2) = \sum_{n=1}^\infty 1/n^2 = \pi^2/6$, which is also as known as the Riemann function at 2 or the Basel problem. Also, the report will figure out important theorems and proofs needed while proving the identity. ~\\


수학세미나I 수업 시간에 $\sum_{n=1}^\infty 1/n^2$의 값을 구하는 토의 학습을 했다. 수업에서는 $f(x)=x$의 푸리에 급수를 이용해 값을 구할 수 있었다. 선생님께서는 다른 획기적인 풀이들도 많다고 하셨기에 찾아봤고, 대부분의 풀이들이 적분이나 테일러 급수와 같은 미적분학 이론을 사용한다는 것을 알게 되었다. 이후 미적분학I 수업을 들으며 얻은 지식들을 토대로 위 문제의 답을 구하는 여러 방법을 제시하고자 한다.

본 보고서는 리만 제타함수, 그리고 p급수의 수렴판정과 깊은 관련을 가진 항등식이자 바젤 문제로도 잘 알려진 식인 $\zeta(2) = \sum_{n=1}^\infty 1/n^2 = \pi^2/6$을 증명하는 다양한 방법에 관해 기술하고 있다. 또한, 항등식을 보이는 과정에서 중요한 정리 및 증명들을 짚고 넘어갈 것이다.~\\\\
\end{abstract}
\begin{dedication2}
{\large{19007 권준우}}\\[5mm]
2021년 미적분학I 5반
\end{dedication2}
\begin{copyright2}
	
	\begin{figure}[h]
		\includegraphics[width=3in]{logo_full}
	\end{figure}
	
	\noindent Copyright \textcopyright\the\year\ Junu Kwon \\ % Copyright of the form
	
	\noindent \textsc{https://www.gs.hs.kr/}\\ % URL
	
	\noindent \textit{Compiled \today} 
	\if@openright\cleardoublepage\else\clearpage\fi
\end{copyright2}\if@openright\cleardoublepage\else\clearpage\fi
    %\tableofcontents*\if@openright\cleardoublepage\else\clearpage\fi
    %\listoffigures*\if@openright\cleardoublepage\else\clearpage\fi
    %\listoftables*\if@openright\cleardoublepage\else\clearpage\fi
    %%% will only print what is used ... useful.
%% also acronyms are clickable, which is awesome

\chapter*{List of Abbreviations}
\markboth{List of Abbreviations}{List of Abbreviations}
               
\begin{acronym}\itemsep-20pt\parsep-20pt %% if you remove these spacing params this list becomes huge!
\acro{GSHS}{Gyeonggi Science High School}
\acro{GPU}{Graphic Processing Unit}
\acro{CPU}{Central Processing Unit}
\acro{GUI}{Graphic User Interface}
\acro{CLI}{Command Line Interface}
\acro{CUDA}{Compute Unified Device Architecture}
\acro{CP}{Compute Compatibility}
\acro{UA}{Used Acronym}
\acro{lox}[\ensuremath{LOX}]{Liquid Oxygen}
\acro{lh2}[\ensuremath{LH_2}]{Liquid Hydrogen}
\acro{IC}{Integrated Circuit}
\acro{BUT}{Block Under Test}
\acrodefplural{BUT}{Blocks Under Test}    
\end{acronym}
\if@openright\cleardoublepage\else\clearpage\fi

%% Note: always use \input as you cannot nest \includes (amongst other things)
\definecolor{dkgreen}{rgb}{0,0.6,0}
\definecolor{gray}{rgb}{0.5,0.5,0.5}
\definecolor{mauve}{rgb}{0.58,0,0.82}
\makeatletter
\newcommand*\verbfont{%
    \normalfont\ttfamily\hyphenchar\font\m@ne\@noligs}
\makeatother
\lstset{ %
  language=bash,                % the language of the code
  basicstyle   = \color{black}\raggedright\verbfont\footnotesize\selectfont,
  numberstyle=\footnotesize,          % the size of the fonts that are used for the line-numbers
  stepnumber=2,                   % the step between two line-numbers. If it's 1, each line 
                                  % will be numbered
  numbersep=5pt,                  % how far the line-numbers are from the code
  backgroundcolor=\color{white},      % choose the background color. You must add \usepackage{color}
  showspaces=false,               % show spaces adding particular underscores
  showstringspaces=false,         % underline spaces within strings
  showtabs=false,                 % show tabs within strings adding particular underscores
  frame=single,                   % adds a frame around the code
  tabsize=2,                      % sets default tabsize to 2 spaces
  captionpos=b,                   % sets the caption-position to bottom
  breaklines=true,                % sets automatic line breaking
  breakatwhitespace=false,        % sets if automatic breaks should only happen at whitespace
  title=\lstname,                   % show the filename of files included with \lstinputlisting;
                                  % also try caption instead of title
  numberstyle=\tiny\color{gray},        % line number style
  keywordstyle=\color{blue},          % keyword style
  commentstyle=\color{dkgreen},       % comment style
  stringstyle=\color{mauve},         % string literal style
  escapeinside={\%*}{*)},            % if you want to add a comment within your code
  morekeywords={*,...}               % if you want to add more keywords to the set
  basicstyle=\ttfamily,
  columns=flexible
}
\pagestyle{gspage}
\mainmatter 
    
\chapter*{Various Methods for \\Calculating $\boldsymbol\zeta$(2)}
\graphicspath{{./chap1/images/}}   
%\section{Introduction}
The Riemann zeta function or Euler–Riemann zeta function, ζ(s), is a mathematical function of a complex variable s, and can be expressed as:
\[
    \zeta(s) = \sum_{n=1}^\infty n^{-s} = \frac{1}{1^s} + \frac{1}{2^s} + \frac{1}{3^s} + \cdots, \textrm{If Re(s) > 1}
\]
And the value of the zeta function at $s = 2$ is a famous problem in mathematical analysis with relevance to number theory, first posed by Pietro Mengoli in 1650. Examining$\zeta(2)$, also knows as the Basel problem asks for the precise summation of the reciprocals of the squares of the natural numbers.
\[
    \sum_{n=1}^\infty \frac{1}{n^2} = \frac{1}{1^2} + \frac{1}{2^2} + \frac{1}{3^2} + \cdots
\]
The problem was first solved by Leonhard Euler in 1734 and the value is well known:  $\frac{\pi^2}{6}$
~\\





Before we start the journey, we need to note some alternative forms of the problem, for ease of proving. The identity we need to prove is
\begin{align}
    \zeta(2) = \sum_{n=1}^\infty \frac{1}{n^2} = \frac{\pi^2}{6} \tag{1}\label{eq:1}
\end{align}~\\
Since it is clear that
\[
    \frac{3}{4}\zeta(2) = \sum_{n=1}^\infty \frac{1}{n^2} -  \sum_{n=1}^\infty \frac{1}{(2n)^2}= \sum_{n=0}^\infty \frac{1}{(2n+1)^2} 
\]~\\
\eqref{eq:1} is equivalent with
\begin{align}
    \sum_{n=0}^\infty \frac{1}{(2n+1)^2} = \frac{\pi^2}{8} \tag{2}\label{eq:2}
\end{align}~\\


So, let's start the journey.\\
~\\
%%----------------
\section{Euler's proof for the Basel problem}
Before we start, let's check out Euler's proof in 1735. Using the Taylor expansion,
\[
    \frac{\sin x}{x} = 1 - \frac{x^2}{3!} + \frac{x^4}{5!}+\cdots
\]
Let the roots of the equation $1 - \frac{x^3}{3!} + \frac{x^5}{5!}+\cdots=0$ be $\alpha_1, \alpha_2, \alpha_3, \cdots$.\\
then $\alpha_1^2, \alpha_2^2, \alpha_3^2, \cdots$ are the roots of $1 - \frac{x}{3!} + \frac{x^2}{5!}+\cdots=0$.~\\

Letting
\[
    p(x)=1 - \frac{x^2}{3!} + \frac{x^4}{5!}+\cdots
\]
$p(x)$ can be also expressed as
\[
    p(x)=(x-\alpha_1^2)(x-\alpha_2^2)(x-\alpha_3^2)\cdots
\]~\\
And comparing the coefficient of $x$ in both expressions,
\[
    \frac{1}{6}=\frac{1}{\alpha_1^2}+\frac{1}{\alpha_2^2}+\frac{1}{\alpha_3^2}+\cdots
\]~\\

Meanwhile, the roots of $\frac{\sin x}{x} = 0$ is $n\pi (n \in \mathbb{N})$ so all $n^2\pi^2$s are the root of $p(x)=0$. Substituting it,
\[
    \sum_{n=1}^\infty \frac{1}{n^2} = \frac{\pi^2}{6}
\]
~\\
\begin{remark}
    Actually, the proof has severe errors handling the concept of infinity, which needs to be proven. However, the idea was genuine and the value was correct.
\end{remark}

%%----------------
%%----------------
\newpage
%%----------------
\section{Proof using the Taylor series of $\arcsin x $~}
The Taylor series of the inverse sine function where $|x| \leq 1$ is
\[
    \arcsin x = \sum_{n=0}^\infty \frac{1\cdot3\cdot5\cdots(2n-1)}{2\cdot4\cdot6\cdots2n} \frac{x^{2n+1}}{2n+1}
\]
Let $x = \sin t$ so that
\begin{align}
    t = \sum_{n=0}^\infty \frac{1\cdot3\cdot5\cdots(2n-1)}{2\cdot4\cdot6\cdots2n} \frac{\sin^{2n+1}t}{2n+1}\tag{3}\label{eq:3}
\end{align}
stands for $|t| \leq \frac{\pi}{2}$.~\\\\

Before moving forward, we must prove that \eqref{eq:3} is uniformly convergent.~\\

Let $a_n = \frac{1\cdot3\cdot5\cdots(2n-1)}{2\cdot4\cdot6\cdots2n}$ then
\[
    \frac{1\cdot3\cdot5\cdots(2n-1)}{2\cdot4\cdot6\cdots2n} < \frac{2\cdot4\cdot6\cdots(2n)}{3\cdot5\cdot7\cdots(2n+1)}
\]
\[
    {a_n}^2 < \frac{1}{2n+1}
\]
Therefore,
\[
    a_n < \sqrt{\frac{1}{2n+1}}
\]
~\\

Using the Weierstrass M-test, $\forall n \in \mathbb{N}$,
\[
     |f_n| = |\frac{1\cdot3\cdot5\cdots(2n-1)}{2\cdot4\cdot6\cdots2n} \frac{\sin^{2n+1}x}{2n+1}| \leq \frac{1\cdot3\cdot5\cdots(2n-1)}{2\cdot4\cdot6\cdots2n} \frac{1}{2n+1} = M_n
\]
and the series $\sum_{n=0}^\infty M_n$ converges by the direct comparison test and the p-series test because
\[M_n = a_n\frac{1}{2n+1} < \frac{1}{\sqrt{(2n+1)^3}}\]~\\


So, the terms in the RHS of \eqref{eq:3} are uniformly convergent. We can now integrate each terms independently. ~\\\\


Finally, using the Wallis formula
\[
    \int_{0}^\frac{\pi}{2} \sin^{2n+1}t dt =  \frac{2\cdot4\cdot6\cdots2n}{1\cdot3\cdot5\cdots(2n+1)}
\]
for integrating \eqref{eq:3} from $0$ to $\frac{\pi}{2}$ gives us
\[
     \frac{\pi^2}{8} = \int_{0}^\frac{\pi}{2} t dt = \sum_{n=0}^\infty \frac{1}{(2n+1)^2}
\]
which is equivalent to \eqref{eq:1}, the identity we want to prove.~\\


\begin{remark}
Without proving that the function is uniformly convergent, we cannot integrate each terms inside the right-hand-side of \eqref{eq:3} independently.
\end{remark}
~\\\\
%-------------------------
\section{Proof using the Fourier series of $f(x)=x(1-x)$~}
Take $f(x)=x(1-x)$. Since $f$ is continuous at $[0,1]$ and $f(0)=f(1)$, the Fourier series of $f$ converges to f pointwise.



This gives us
\[
    x(1-x)=\frac{1}{6}-\sum_{n=1}^{\infty}\frac{\cos 2\pi nx}{\pi^2n^2}
\]
and putting $x=0$ we get \eqref{eq:1}.

Alternatively putting $x=\frac{1}{2}$ gives us
\[
    \frac{\pi^2}{12}=-\sum_{n=1}^{\infty}\frac{(-1)^n}{n^2}
\]
which is also equivalent to \eqref{eq:1}.

\newpage
\section{Proof using the Fourier series of $f(x)=x$~}
Lets start with the Fourier expansion of $f(x) = x$ which is:
\[
    a_n =  \frac{1}{\pi}\int_{-\pi}^{\pi}f(x)\cos(nx)dx = 0
\]
\[
    b_n =  \frac{1}{\pi}\int_{-\pi}^{\pi}f(x)\sin(nx)dx = (-1)^{n+1}\frac{2}{n}
\]
~\\
\begin{remark}
    We should prove that the function is uniformly continuous while following the Fourier expansion procedure. However, I will skip it since it's similar to the one we will prove later at the third proof.
\end{remark}
~\\


Using the Parseval's equality
\[
    \frac{1}{\pi}\int_{-\pi}^{\pi}f(x)^2dx = \frac{1}{2}a_{0}^2+\sum_{n=1}^{\infty}(a_n^2+b_n^2)
\]
\[
    \frac{1}{\pi}\int_{-\pi}^{\pi}x^2~dx=\sum_{k=1}^\infty\frac{4}{n^2}
\]~\\


Using 
\[
    \frac{1}{\pi}\int_{-\pi}^{\pi}x^2~dx = \frac{2\pi^2}{3}
\]
We can prove the identity.
\[
    \sum_{n=1}^\infty\frac{1}{n^2}=\frac{\pi^2}{6}
\]~\\
\begin{remark}
    This was the proof which was suggested by the teacher at class. The concept of Fourier expansion is somehow complicated. However, after understanding it, the proof is very easy to follow. Also, we can get the infinite sum $\sum_{n=1}^\infty\frac{1}{n^4}$ using the expansion of $f(x)=x^2$, via the same procedure.
\end{remark}

\newpage
\section{Proof using the Taylor series of $\arctan x$~}
The Taylor series of the inverse tangent function is
\[
    \arctan x = \sum_{n=1}^\infty(-1)^n\frac{x^{2n+1}}{2n+1}
\]
Substituting $x = 1$, we can get the Gregory's formula
\[
    \frac{\pi}{4}=\sum_{n=1}^\infty\frac{(-1)^n}{2n+1} = \frac{1}{1} - \frac{1}{3} + \frac{1}{5} - \cdots
\]
~\\

Rewriting the formula as $\lim_{N\rightarrow\infty}a_N=\frac{\pi}{2}$ where
\[
    a_N = \sum_{n=-N}^{N}\frac{(-1)^n}{2n+1}
\]
Let
\[
    b_N = \sum_{n=-N}^{N}\frac{1}{(2n+1)^2}
\]
So, $\lim_{N\rightarrow\infty}b_N=\frac{\pi^2}{4}$ is consist with \eqref{eq:2}, which is the identity we want to prove. So, we shall show that $\lim_{N\rightarrow\infty}(a_N^2-b_N)=0$~\\

If $n \neq m$ then
\[
    \frac{1}{(2n+1)(2m+1)}=\frac{1}{2(m-n)}(\frac{1}{2n+1} - \frac{1}{2m+1})
\]
and so

\begin{equation}
\begin{split}
a_N^2-b_N~
={}& ~{\sum_{n=-N}^{N}\sum_{m=-N}^{N}}^\prime \frac{(-1)^{m+n}}{2(m-n)}(\frac{1}{2n+1} - \frac{1}{2m+1})\\
={}& ~{\sum_{n=-N}^{N}\sum_{m=-N}^{N}}^\prime \frac{(-1)^{m+n}}{(2n+1)(m-n)}\\
={}& ~\sum_{n=-N}^{N}\frac{(-1)^{n} c_{n,N}}{(2n+1)} \nonumber
\end{split}
\end{equation}
where the prime on the summations means that the terms with zero denominators, where $n=m$ are omitted, and
\[
    c_{n,N}={\sum_{m=-N}^{N}}^\prime \frac{(-1)^{m}}{(m-n)}
\]

It is obvious that $c_{-n,N}=-c_{n,N}$ and $c_{0,N}=0$. For $n > 0$,
\[
    c_{n,N}=(-1)^{n+1}\sum_{j=N-n+1}^{N+n}\frac{(-1)^j}{j}
\]
which means that $|c_{n,N}| \leq \frac{1}{N-n+1}$ since the magnitude of an converging altering sum is smaller than that of the first term. Hence
\begin{equation}
\begin{split}
|a_N^2-b_N| ~\leq{}& ~\sum_{n=1}^N(\frac{1}{(2n-1)(N-n+1)} + \frac{1}{(2n+1)(N-n+1)})\\
={}& ~\sum_{n=1}^N\frac{1}{2N+1}(\frac{2}{(2n-1)} + \frac{1}{(N-n+1)})\\
{}& ~+ \sum_{n=1}^N\frac{1}{2N+3}(\frac{2}{(2n+1)} + \frac{1}{(N-n+1)})\\
\leq{}& ~\frac{1}{2N+1}(2+4\log(2N+1)+2+2log(N+1)) \nonumber
\end{split}
\end{equation}
and so $\lim_{N\rightarrow\infty}(a_N^2-b_N)=0$ as required.~\\\\

\section*{Ending the journey}
We could prove the Basel problem using various methods using the Fourier series of $x$ or the Taylor series of $\arcsin x$ and $\arctan x$. Although it seems enough, there are plenty more proofs for the Basel problem outside the world. I'm afraid that the journey ends here, but however I could realize that the knowledge of math diverges to infinity.~\\


\vspace{4.5cm}
\begin{figure}[H]
\begin{table}[H]
\begin{center}
\begin{tabular}{l}
\hline
\textbf{~~~이 보고서는 2021년 미적분학I 자율보고서로 작성되었습니다.~~~}\\
\hline
\end{tabular}
\end{center}
\end{table}
\end{figure} 
    %
\graphicspath{{./chap2/images/}}
\chapter{Convergence}
\section{The integral test}
\section{Compression tests}
    %
\graphicspath{{./chap3/images/}} 

\chapter{Absolute Convergence}
\section{The ratio test}
\section{The root test}
    %\appendix
        %\chapter{Troubleshooting}
\section{Graphic Driver}
\subsection{nvidia-smi}
nvidia driver labrary mismatch~\\

cuda
\section{Storage}
dev1
\section{Python Environments}
pip, pip3, venv
\section{Jupyter}
config
\section{SSH}
sshd\_config
\section{Database}
mongodb, mariadb
\section{Nginx}
nginx.conf     % these are just test names as I didn't know what you'd want
        %\chapter{GSHS CSLAB}
정보과 연구용 기기 목록
\begin{itemize}
    \item Nvidia RTX 2070 eGPU $\times$ 2
    \item Nvidia RTX 2080 Super eGPU $\times$ 1
    \item Nvidia RTC 3070 $\times$ 1 
    \item Intel NUC $\times$ 2
    \item 본 서버.
\end{itemize}
    
        %
\chapter{Contributors}

19007 권준우, 19011 김도영, 19032 노현서, 19042 박영민
 

{\backmatter
    % Bibliography
    \if@openright\cleardoublepage\else\clearpage\fi
    \bibliographystyle{gs-plainnat} %% specific plainnat does not show url for articles
    {\footnotesize\bibliography{chap1/introduction_biblio,chap2/background_and_lit_overview_biblio}}
	\printindex
}

\end{document}

%%% The End %%%